\documentclass[10pt]{article}
\usepackage[T1]{fontenc}
\usepackage[utf8]{inputenc}
\usepackage[french]{babel}
\usepackage{amsmath,amssymb,mathrsfs,mathpazo}
\usepackage{geometry}
\geometry{a4paper,margin=2cm}
\usepackage{graphicx,array,tabularx,tabularray,colortbl,multirow,multicol,booktabs,enumitem,pifont,ulem,fancybox,diagbox,cancel}
\usepackage{pstricks,pstricks-add,pst-plot,pst-tree,pst-node,pst-text,pst-func,pst-math,pst-bspline,pst-3dplot,pst-eucl,pst-solides3d}
\let\pscalculate\undefined
\usepackage{pst-bezier}
\usepackage{tikz,pgfplots,tkz-tab}
\usetikzlibrary{arrows,arrows.meta,patterns,calc,positioning,trees,decorations.markings,shapes.geometric,shapes.misc,shapes.arrows,intersections,decorations.pathmorphing,decorations.pathreplacing,backgrounds,shadows,babel}
\pgfplotsset{compat=1.18}
\usepgfplotslibrary{fillbetween}
\usepackage[np]{numprint}
\usepackage{listings,makecell,pas-tableur,twoopt}
\let\barre\undefined
\usepackage[pstricks]{tablvar}
\definecolor{darkgreen}{RGB}{0,100,0}
\providecommand{\red}{\color{red}}
\providecommand{\blue}{\color{blue}}
\providecommand{\green}{\color{green}}
\providecommand{\black}{\color{black}}
\providecommand{\white}{\color{white}}
\providecommand{\gray}{\color{gray}}
\providecommand{\cyan}{\color{cyan}}
\providecommand{\magenta}{\color{magenta}}
\providecommand{\yellow}{\color{yellow}}
\providecommand{\violet}{\color{violet}}
\providecommand{\orange}{\color{orange}}
\providecommand{\euro}{\texteuro}
\providecommand{\EUR}{\texteuro}
\providecommand{\ieme}{\textsuperscript{e}}
\providecommand{\degres}{\ensuremath{^\circ}}
\providecommand{\degree}{\ensuremath{^\circ}}
\providecommand{\vv}[1]{\overrightarrow{#1}}
\providecommand{\Coord}[1]{#1}
\providecommand{\decosix}{}
\providecommand{\starredbullet}{\ensuremath{\bullet}}
\pagestyle{empty}
\setlength{\parindent}{0pt}
\usepackage{adjustbox}\DeclareUnicodeCharacter{2236}{:}\DeclareUnicodeCharacter{2190}{\ensuremath{\leftarrow}}
\providecommand{\R}{\mathbb{R}}
\providecommand{\N}{\mathbb{N}}
\providecommand{\Z}{\mathbb{Z}}
\providecommand{\Q}{\mathbb{Q}}
\providecommand{\D}{\mathbb{D}}
\providecommand{\C}{\mathbb{C}}
\providecommand{\e}{{\mathrm{e}}}
\providecommand{\vect}[1]{\overrightarrow{#1}}
\providecommand{\vectt}[1]{\overrightarrow{\mathrm{#1}}}
\providecommand{\vvt}[1]{\overrightarrow{\mathrm{#1}}}
\providecommand{\np}[2][]{{#2~#1}}
\providecommand{\vv}[1]{{\overrightarrow{#1}}}
\renewcommand{\d}{{\mathrm{d}}}
\providecommand{\degres}{{^\circ}}
\providecommand{\degre}{{^\circ}}
\providecommand{\blue}{\color{blue}}
\providecommand{\red}{\color{red}}
\providecommand{\enskip}{\quad}
\providecommand{\llbracket}{\lbrack\!\lbrack}
\providecommand{\rrbracket}{\rbrack\!\rbrack}
\providecommand{\euro}{{\text{€}}}
\providecommand{\no}{{\text{n°}}}
\providecommand{\interval}[2]{{[#1;#2]}}
\providecommand{\intervFO}[2]{{[#1;#2[}}
\providecommand{\dotfill}{\ldots}

\providecommand{\ds}{\displaystyle}
\providecommand{\textding}[1]{#1}
\providecommand{\Oij}{\(\mathrm{O};\vec{\imath},\vec{\jmath}\)}
\providecommand{\Oijk}{\(\mathrm{O};\vec{\imath},\vec{\jmath},\vec{k}\)}
\providecommand{\Ouv}{\(\mathrm{O};\vec{u},\vec{v}\)}
\providecommand{\intervFF}[2]{\ensuremath{[#1;#2]}}
\renewcommand{\euro}{\ifmmode\text{\texteuro}\else\texteuro\fi}
\DeclareUnicodeCharacter{00B0}{\ifmmode^{\circ}\else\textdegree\fi}
\DeclareUnicodeCharacter{20AC}{\ifmmode\text{\texteuro}\else\texteuro\fi}
% Énoncés : Éducation nationale. Transcription et corrigés : APMEP, Denis Vergès et contributeurs. Découpage et classement : CoopMaths, Sébastien Lozano, Éric Elter et contributeurs. Adaptation de la présentation : Maths Explorer.
% CC BY-SA : https://coopmaths.fr/www/about/
% https://coopmaths.fr/alea/static/bac/2021/tex/bac_2021_05_ameriquenord_3_cor.tex
\begin{document}
\noindent\textit{APMEP / CoopMaths. CC BY-SA. Présentation adaptée par Maths Explorer.}\par\medskip
\textbf{Commun à tous les candidats}

\smallskip








\begin{center}
\psset{xunit=0.85cm,yunit=0.9cm}
\begin{pspicture}(6,7)
\pspolygon(0.5,0.4)(5.5,0)(5.5,5)(0.5,5.4)
\uput[dl](0.5,0.4){A} \uput[dr](5.5,0){B} \uput[u](5.5,5){F} \uput[ul](0.5,5.4){E}
\psline(5.5,0)(8.5,1.4)(8.5,6.4)(5.5,5)
\uput[r](8.5,1.4){C} \uput[ur](8.5,6.4){G} 
\psline(8.5,6.4)(3.5,6.8)(0.5,5.4)
\uput[u](3.5,6.8){H} \uput[u](3,5.2){I}\uput[dr](7,0.7){J}\uput[l](0.5,2.9){K}
\psline[linewidth=1.6pt](0.5,0.4)(3,5.2)
\psline[linestyle=dashed,linewidth=1.6pt](3,5.2)(7,0.7)
\psline[linestyle=dashed,linewidth=1.6pt](0.5,2.9)(3.5,6.8)
\psline[linestyle=dashed](0.5,0.4)(3.5,1.8)(3.5,6.8)
\uput[ur](3.5,1.8){D}
\psline[linestyle=dashed](3.5,1.8)(8.5,1.4)
\end{pspicture}
\end{center}

En prenant le repère orthonormé $\left(\text{A}~;~\vect{\text{AB}},~ \vect{\text{AD}},~ \vect{\text{AE}}\right)$ on a :

\[\text{A}(0~;~0~;~0),\, \text{B}(0~;~0~;~0),\, \text{C}(1~;~1~;~0),\, \text{D}(0~;~1~;~0),\, \text{E}(0~;~0~;~1)\]

\[\, \text{F}(1~;~0~;~1),\, \text{G}(1~;~1~;~1),\, \text{H}(0~;~1~;~1),\, \text{I}(0,5~;~0~;~1),\, \text{J}(1~;~0,5~;~0),\, \text{K}(0~;~0~;~0,5)\]
\begin{enumerate}
\item 

On a $\vect{\text{AI}}(0,5~;~0~;~1)$ et $\vect{\text{KH}}(0~;~1~;~0,5)$ : ces vecteurs ne sont pas colinéaires, donc les droites (AI) et (KH) ne sont pas parallèles.


\item 
	\begin{enumerate}
		\item 
Voir plus haut.
		\item 
On a $\vect{\text{IJ}}(0,5~;~0,5~;~-1), \, \vect{\text{AE}}(0~;~0~;~1)\, \vect{\text{AC}}(1~;~1~;~0)$.

On a $2\vect{\text{IJ}} + 2\vect{\text{AE}} = \vect{\text{AC}}$.

Le vecteur $\vect{\text{AC}}$ est donc une combinaison des vecteurs $\vect{\text{IJ}}$ et $\vect{\text{AE}}$ : ces trois vecteurs sont donc coplanaires.
	\end{enumerate}

	














\item 

$d_1$ a pour vecteur directeur $\vect{u_1}(1~;~-2~;~3)$ et $d_2$ a pour vecteur directeur $\vect{u_2}(1~;~1~;~2)$ : ces vecteurs ne sont pas colinéaires, donc les droites $d_1$ et $d_2$ ne sont pas parallèles 
\item 

Le plan a pour vecteur normal le vecteur $\vect{p}(1~;~3~;~-2)$ et $d_2$ a pour vecteur directeur $\vect{u_2}(1~;~1~;~2)$.

Or $\vect{p} \cdot \vect{u_2} =  1 + 3 - 4 = 0$ : les vecteurs sont orthogonaux donc la droite $d_2$ est parallèle au plan $\mathcal P$.
\item 
\textbf{Méthode 1}

Soit $\Delta$ la perpendiculaire à $\mathcal P$ contenant M. Cette droite a pour vecteur directeur le vecteur $\vect{p}$, donc une équation paramétrique de $\Delta$ est :

$\left\{\begin{array}{l c l}
x	&=&5 + 1t\\
y 	&=&3 + 3t\\
z	&=&1 - 2t\\
\end{array}\right., t \in \R.$

Le projeté L,  de M sur le plan $\mathcal P$ a ses coordonnées qui vérifient les quatre équations :

$\left\{\begin{array}{l c l}
x	&=&5 + 1t\\
y 	&=&3 + 3t\\
z	&=&1 - 2t\\
x + 3y - 2z + 2&=&0
\end{array}\right. , t \in \R. \Rightarrow 5 + t  + 3(3 + 3t) - 2(1 - 2t) + 2 = 0 \iff$

$5 + t + 9 + 9t - 2 + 4t + 2 = 0 \iff 14t  + 14 = 0 \iff t + 1 = 0 \iff t = - 1$.

En reportant dans les trois premières équations du système, on trouve les coordonnées de L projeté orthogonal de M sur $\mathcal P$ :

$\left\{\begin{array}{l c l}
x	&=&5 - 1\\
y 	&=&3 + 3\times (- 1)\\
z	&=&1 - 2\times (- 1)\\
\end{array}\right. \iff \left\{\begin{array}{l c l}
x	&=&4\\
y 	&=&0\\
z	&=&3\\
\end{array}\right.$

Donc le projeté orthogonal de M sur le plan $\mathcal P$ est le le point L(4~;~0~;~3).
\end{enumerate}

\textbf{Méthode 2}

On a $\vect{\text{ML}}(-1~;~- 3~;~2)$, donc $\vect{\text{ML}} = - \vect{p}$ est un vecteur normal au plan $\mathcal{P}$.

D'autre part L(4~;~0~;~3) $\in \mathcal{P} \iff 4 + 3 \times 0 - 2 \times 3 + 2 = 6 - 6 = 0$ est vraie, donc L est le projeté orthogonal de M sur le plan $\mathcal P$.
\bigskip

\end{document}
